\subsection{Resumen}
\label{subsec:resumen}

% ####################################################################################################################
% ####################################################################################################################

\todo{P1 -- Ley de Ohm: valor de $R$ obtenido por regresión lineal, pendiente $m$, ordenada $b$, verificación de linealidad.}

% ####################################################################################################################
% ####################################################################################################################

\textbf{Práctica 2 -- Medición de resistencias:} Se midieron tres resistencias con un multímetro digital (UT58D) en modo ohmímetro, estimando la incerteza a partir de las especificaciones del fabricante. Los resultados obtenidos fueron:

\begin{itemize}
    \item $R_1 = (5{,}6 \pm 0{,}1)\ \text{M}\Omega$
    \item $R_2 = (562 \pm 8)\ \text{k}\Omega$
    \item $R_3 = (5{,}50 \pm 0{,}05)\ \text{k}\Omega$
\end{itemize}

Los tres valores resultaron compatibles con el valor nominal dentro de la tolerancia declarada por el fabricante ($\pm 5\%$).

% ####################################################################################################################
% ####################################################################################################################

\todo{P3 -- Influencia del aparato: verificación de la 2ª Ley de Kirchhoff para cada par de resistencias, efecto de $R_{int} = 10\ \text{M}\Omega$ del tester sobre la lectura.}
